% The constraint-promotion barrier in perturbative reduction of higher-derivative
% Hamiltonians.  Documents an architectural limitation hit during TIDAL's v6
% Phase 2 (Lagrangian Perturbative Substitution, LPS); cites established
% physics analogues (vDVZ, Boulware-Deser, Stuckelberg) and surveys methods
% in the literature that fail for this case.
%
% Uses macros from preamble.tex.

\section{The Constraint-Promotion Barrier in Perturbative Hamiltonian Reduction}
\label{sec:pr-constraint-barrier}

\subsection{Abstract}
\label{sec:pr-cb-abstract}

We document an architectural limitation of perturbative reduction methods for
higher-derivative Hamiltonians, hit during TIDAL's v6 Phase~2 (Lagrangian
Perturbative Substitution, LPS).  When a small parameter~$\varepsilon$ promotes
an algebraic-constraint field at $\varepsilon = 0$ to a higher-derivative dynamical
field at $\varepsilon \ne 0$, the constraint structure changes \emph{discontinuously}
across the critical surface.  Substitution-based perturbative-reduction methods
in the literature (Jaen--Llosa--Molina~1986 and descendants; effective-field-theory
field redefinitions; Cayuso--Lehner~2017) all share an architectural assumption
that excludes this case.  Faddeev--Jackiw symplectic reduction, the cubic-PGT
redesign of Bahamonde \emph{et al.}~2025, and DHOST degeneracy conditions
(Langlois~2018) likewise do not provide a perturbative bridge.  TIDAL's
Phase~2 LPS therefore aborts with a clear diagnostic when the constraint-promotion
case is detected (issue~\#321), and the affected theories
($b_5\,\tilde{R}^2$ in \texttt{graviton\_torsion}, \texttt{torsion\_gertsenshtein},
\texttt{torsion\_gertsenshtein\_combined}) retain their pre-LPS Hamiltonian
JSONs.  The barrier is structurally analogous to the van~Dam--Veltman--Zakharov
(vDVZ) discontinuity in massive gravity, the Boulware--Deser non-decoupling, and
the Maxwell--Proca constraint reclassification at $m = 0$.  Three research
directions are identified for future work; none has a published implementation
for the case at hand.

\subsection{Introduction}
\label{sec:pr-cb-intro}

The TIDAL pipeline (\cref{sec:perturbative-reduction}) reduces higher-derivative
Lagrangians to physically meaningful second-order evolution via Parker--Simon
iterative reduction~\cite{ParkerSimon1993}, generalised in v6 to multi-field
linearised PDEs with kinetic-coefficient promotion (Gap~B/C of the v6 plan;
\cref{sec:pr-v6}).  This iterative scheme is theory-agnostic and ghost-free
\emph{at the level of the equations of motion}: the LHS operator at every
order is the second-order base operator, and higher-derivative terms appear
only as sources to a Duhamel kernel evaluated on the previous-order solution.

Energy measurement and Noether consistency, however, require a Hamiltonian.
TIDAL's v6 Phase~D extracts the canonical Hamiltonian via the standard
Legendre transform of the (Phase~B IBP-reduced) Lagrangian.  For perturbative
theories with $\partial_t^2$ residues that IBP cannot eliminate
(\cref{sec:pr-cb-residue}), the v6 Phase~2 Lagrangian Perturbative Substitution
(LPS, \cref{sec:pr-cb-lps}) was introduced to apply the order-0 EOM as a
substitution rule on the Lagrangian, thereby producing a fully second-order
input to the Legendre transform.

This paper documents an architectural limitation of LPS encountered during
its initial deployment.  For the family of theories central to TIDAL's
target science --- $b_5\,\tilde{R}^2$ in Poincar\'e Gauge Theory, where
$b_5$ is a small coupling --- the small parameter promotes algebraic-constraint
fields at $b_5 = 0$ to fourth-order dynamical fields at $b_5 \ne 0$.  The
substitution-based LPS algorithm cannot bridge this discontinuity because
there is no $\partial_t^2$~term in the order-0 EOM to solve for.  A
literature survey (Section~\ref{sec:pr-cb-methods}) confirms that every
published reduction technique --- substitution, field redefinition,
Faddeev--Jackiw, fixing-the-equations, DHOST, cubic redesign --- shares the
same architectural assumption and fails in the same way.

The architectural barrier is structurally analogous to several
well-established discontinuous limits in field theory
(Section~\ref{sec:pr-cb-analogues}): the vDVZ discontinuity in massive
gravity, the Boulware--Deser non-decoupling, the Maxwell--Proca constraint
reclassification.  Casting TIDAL's barrier in this established language
clarifies what is and is not novel.

The remainder of this section presents the architectural analysis
(Sections~\ref{sec:pr-cb-residue}--\ref{sec:pr-cb-lps}), the discontinuous-limit
framing (Section~\ref{sec:pr-cb-analogues}), the method survey
(Section~\ref{sec:pr-cb-methods}), TIDAL's resolution
(Section~\ref{sec:pr-cb-resolution}), and the open research directions
(Section~\ref{sec:pr-cb-future}).

\subsection{Higher-derivative residues and integration by parts}
\label{sec:pr-cb-residue}

A linearised Lagrangian containing the curvature-squared term $\tilde{R}^2$
expands, after component decomposition (Phase~A), to a polynomial in the
metric perturbation components $h_{\mu\nu}$ and their partial derivatives.
Among the resulting terms are products of the form
\begin{equation}
  \mathcal{L}_1 \supset b_5 \, (\partial_t^2 h_a)(\partial_t^2 h_b) \,,
  \label{eq:pr-cb-ddtsqsq}
\end{equation}
which contain pure second time derivatives of the dynamical (or
constraint-promoted) components of the metric perturbation.

The standard Phase~B integration-by-parts (IBP) on the time variable
\cite{ParkerSimon1993} converts $f \cdot \partial_t^2 g \to -(\partial_t f)\,(\partial_t g)$
when one of the factors carries no time derivative.  IBP successfully
reduces $f \cdot \partial_t^2 g$ products to first-order time derivatives.
But it \emph{cannot} reduce
\eqref{eq:pr-cb-ddtsqsq}: every IBP move on $(\partial_t^2 h_a)(\partial_t^2 h_b)$
either leaves the time-derivative order at two (passing the derivative
between factors, $(\partial_t^2 h_a)(\partial_t^2 h_b) \to -(\partial_t^3 h_a)(\partial_t h_b)$,
which \emph{raises} the maximum order) or fails at the boundary terms.
The irreducible $\partial_t^2$ residue in \eqref{eq:pr-cb-ddtsqsq} is
fundamental to the algebraic-IBP path.

In TIDAL the irreducible residues manifest as JSON Hamiltonian operators
of the form \texttt{mixed\_2\_0\_0\_0} (encoding $\partial_t^2$ in 1+1D
spatial indices).  Direct inspection of the affected examples
(\cref{sec:pr-cb-resolution}) confirms 3 such operators in
\texttt{graviton\_torsion} and \texttt{torsion\_gertsenshtein}, and 33 in
\texttt{torsion\_gertsenshtein\_combined}, all carrying coefficient
proportional to the small parameter ($b_5$ or one of $a_1,\dots,a_6$ in
the combined theory).

\subsection{Lagrangian Perturbative Substitution (LPS)}
\label{sec:pr-cb-lps}

To eliminate the \texttt{mixed\_2\_0\_0\_0} residues at the Lagrangian
level (so the downstream Legendre transform operates on a strictly
second-order input), v6 Phase~2 introduces the Lagrangian Perturbative
Substitution module
(\href{https://github.com/WilliamRoyce/torsion-gertsenshtein/blob/main/tidal/wolfram/PerturbativeReduction.wl}{\texttt{tidal/wolfram/PerturbativeReduction.wl}},
commit \texttt{742c01a}).  LPS implements the algorithm:

\begin{enumerate}
  \item \emph{Extract base equations}: set every declared small parameter to
        zero in each Phase~B Euler--Lagrange equation, yielding the order-0
        EOM $E_{\mathrm{base}}[q] = 0$.
  \item \emph{Solve for $\partial_t^2$}: solve the order-0 EOMs simultaneously
        for $\partial_t^2 q_i$ for each dynamical field, producing the rule
        $\partial_t^2 q_i \to f_i(q, \partial_t q, \partial_x q, \dots)$.
  \item \emph{Chain higher orders}: build rules for $\partial_t^3 q_i,
        \partial_t^4 q_i, \dots$ by differentiating $f_i$ with respect
        to~$t$ and substituting the lower-order rules.
  \item \emph{Apply and truncate}: apply the rules to the Lagrangian via
        repeated replacement (\verb|//.|) and truncate the result at the
        declared perturbative order via Series.
\end{enumerate}

For theories where every field with $\partial_t^2$ in the Lagrangian is also
\emph{dynamical at order zero} --- the standard case described by
Parker--Simon and the EFT-reduction literature --- this algorithm completes
cleanly.  Pais--Uhlenbeck-style perturbations to dynamical Lagrangians
satisfy this case~\cite{BenderMannheim2008,Woodard2015}, as does
Euler--Heisenberg in TIDAL's representation (whose component-decomposed
Lagrangian is naturally first-order in time after IBP).

The algorithm \emph{fails} when the Lagrangian contains
$\partial_t^2$~terms of a field~$h_c$ that has \emph{no} $\partial_t^2 h_c$
term in its order-0 EOM~--- because Step~2 has no rule to construct.
This is the constraint-promotion case.

\subsection{The constraint-promotion case}
\label{sec:pr-cb-promotion}

\paragraph{Definition.}
A theory exhibits \emph{constraint promotion} if a small parameter $\varepsilon$
in the Lagrangian (with $\varepsilon = 0$ giving a healthy second-order base
theory) raises the time-derivative order of one or more component fields'
Euler--Lagrange equations.  Equivalently: the kinetic coefficient
$M(\varepsilon) = M_0 + \varepsilon M_1 + \dots$ for the affected component(s)
satisfies $M_0 = 0$, so at $\varepsilon = 0$ the time-derivative term vanishes
and the equation becomes algebraic; for $\varepsilon \ne 0$ the leading
$\varepsilon M_1$ term restores time evolution at higher order.

\paragraph{Concrete example: $b_5\,\tilde{R}^2$ in PGT.}
Let $\tilde{R}^{abc} \equiv \frac{1}{2}\,\varepsilon^{abcd}\, R_{de}{}^{de}$ be the
parity-odd dual of the curvature.  The Lagrangian contribution
$\mathcal{L}_1 = b_5\,\tilde{R}_{abc}\tilde{R}^{abc}$ produces, after component
decomposition for the metric perturbation $h_{\mu\nu}$ with TT gauge in
1+1D, a kinetic structure
\begin{equation}
  M(b_5)\,\partial_t^2 h_a = -K_a[h, \partial_x h, \dots] \,, \qquad
  M_a(b_5) = c_a\, b_5 \quad\text{for } a \in \{4, 7, 9\}.
  \label{eq:pr-cb-Mb5}
\end{equation}
The vanishing of $M_a(0) = 0$ for $a \in \{4, 7, 9\}$ is the constraint-promotion
signature.  At $b_5 = 0$, components $h_4, h_7, h_9$ satisfy algebraic
constraints; at $b_5 \ne 0$, they become \emph{fourth}-order dynamical fields
(after Euler--Lagrange variation of $\tilde{R}^2$ and component decomposition).
Direct inspection of TIDAL's v6 spec parser
(\texttt{tidal/symbolic/json\_loader.py::canonicalize\_kinetic\_for\_perturbation})
classifies these three fields exactly as ``constraint at base, dynamical at
correction''.  The Phase~2 LPS detects this case via the absence of
$\partial_t^2 h_a$ in the order-0 EOM and \emph{aborts}.

\paragraph{Why algebraic substitution fails.}
Step~2 of the LPS algorithm (\cref{sec:pr-cb-lps}) requires solving
$E_{\mathrm{base}}[q] = 0$ for $\partial_t^2 q_i$.  For a constraint-promoted
field $h_c$, the order-0 EOM has the form
\begin{equation}
  E_{\mathrm{base}}[h_c] = \alpha_c\, h_c + \beta_c[\,\text{other fields and their derivatives}\,]
  \label{eq:pr-cb-base-algebraic}
\end{equation}
with no $\partial_t^2 h_c$ term.  The substitution rule
$\partial_t^2 h_c \to (\dots)$ does not exist; one can only solve
\eqref{eq:pr-cb-base-algebraic} for $h_c$ itself, in terms of the other
fields:
\begin{equation}
  h_c \to -\beta_c / \alpha_c \,.
\end{equation}
Substituting this into the Lagrangian's $b_5\,(\partial_t^2 h_c)^2$ term
introduces $\partial_t^2(\beta_c/\alpha_c)$, which by the chain rule becomes
$\partial_t^2$ of \emph{other} fields --- which themselves must then be
substituted using \emph{their} order-0 rules.  In a multi-field gauge
theory the chain rule generates rapidly proliferating substitutions that
were observed in practice (commit \texttt{742c01a}, branch attempt) to
fail to converge within reasonable resources ($> 5$~min,~$> 2$~GB on the
38-field \texttt{torsion\_gertsenshtein} theory).  More fundamentally, the
chain rule converts $b_5\,(\partial_t^2 h_c)^2$ into terms that are
\emph{fourth}-order in time derivatives of other dynamical fields ---
exactly the obstruction LPS was meant to remove.  The barrier is
not a performance issue but a structural one.

\subsection{Established physics analogues: discontinuous limits}
\label{sec:pr-cb-analogues}

The mathematical structure --- a parameter-dependent constraint count that
changes discontinuously across $\varepsilon = 0$ --- is well known in
field theory under several names~\cite{deRham2014,Hinterbichler2012}.
Following the conventions of those reviews, we adopt the term
``\emph{discontinuous limit of the constraint structure}'' (or
``singular limit'') and avoid the thermodynamics-flavoured ``phase
transition'' (no order parameter is present).  This subsection presents
the established analogues.

\paragraph{The vDVZ discontinuity (massive gravity).}
The smooth $m \to 0$ limit of Pauli--Fierz massive
gravity~\cite{vanDamVeltman1970,Zakharov1970} does \emph{not} recover
linearised general relativity: there is a finite mismatch in the
predicted bending of light around the Sun (factor $4/3$).  Mechanically:
massive spin-2 has 5 propagating degrees of freedom; massless spin-2 has
2.  In the $m \to 0$ limit the helicity-2 (2~dof), helicity-1 (2~dof),
and helicity-0 (1~dof) modes do \emph{not all decouple} --- the helicity-0
scalar retains finite coupling to the trace of the stress-energy
tensor~\cite[\S III]{Hinterbichler2012}.  The Stuckelberg trick exposes
this systematically: auxiliary fields propagate at $m \ne 0$ but become
pure-gauge constraints at $m = 0$.  The structural change is identical to
TIDAL's case: $\varepsilon$ promotes a constraint to a dynamical degree of
freedom; the limit $\varepsilon \to 0$ is non-smooth.

\paragraph{The Boulware--Deser non-decoupling.}
The 1972 paper~\cite{BoulwareDeser1972} demonstrated that the
Pauli--Fierz tuning that gives 5~dof at quadratic order is \emph{unstable}
under generic nonlinear extension: a 6th (ghost) mode appears off the
Pauli--Fierz surface in coefficient space.  This is structurally identical
to TIDAL's situation: at one fine-tuned point in coupling space the
constraint count is one thing; perturbing the coefficients adds a
propagating mode.  The dRGT~\cite{deRham2014} resolution requires
fine-tuning of \emph{multiple} coefficients to preserve the constraint
non-perturbatively --- not a generic perturbative reduction.

\paragraph{The Maxwell--Proca boundary.}
The Stuckelberg field~\cite{RueggRuizAltaba2003} provides the cleanest
toy example: Maxwell electrodynamics has 2~first-class constraints
(Gauss's law and tracelessness of $A_0$), giving 2~dof.  Proca theory
($+\frac{1}{2}m^2 A_\mu A^\mu$) has 2~second-class constraints, giving
3~dof.  At $m = 0$ first-class becomes second-class, and the constraint
that determined $A_0$ becomes a dynamical equation.  This is the textbook
prototype of the constraint reclassification at the boundary.

\paragraph{The mathematical underpinning.}
Henneaux \& Teitelboim~\cite[\S 1.4.3]{HenneauxTeitelboim1992} formalise
this as an ``\emph{irregular constraint surface}'' --- a point in
coupling-parameter space where the rank of the constraint Jacobian
drops.  The canonical-analysis literature acknowledges this obstruction
but does \emph{not} provide a perturbative method that crosses it.  A
modern Faddeev--Jackiw treatment~\cite{DiazSaldivia2025} reads the same
phenomenon off the rank of the symplectic 2-form: at the critical surface
the rank jumps, signalling the appearance or disappearance of constraints.

\paragraph{Summary.}
The constraint-promotion case is not new physics; it is a particular
instance of a class of \emph{discontinuous (singular) limits of the
constraint structure} that includes vDVZ, Boulware--Deser non-decoupling,
and Stuckelberg.  In every case the perturbative expansion in
$\varepsilon$ around $\varepsilon = 0$ does not converge to the
$\varepsilon = 0$ theory at any finite order, because the kinematic
phase space changes discontinuously at the boundary.

\subsection{Survey of perturbative-reduction methods and why they fail}
\label{sec:pr-cb-methods}

We have surveyed the literature on perturbative reduction of higher-derivative
Lagrangians.  Each major method is discussed below, with explicit attention
to the constraint-promotion case.  The unifying conclusion: every method
shares an architectural assumption --- that fields propagating at
$\varepsilon \ne 0$ are also propagating at $\varepsilon = 0$ --- and is
therefore structurally unable to handle constraint promotion.

\paragraph{Jaen--Llosa--Molina (JLM) and descendants.}
The original on-shell-reduction method~\cite{JLM1986} substitutes
$\partial_t^2$~terms in the higher-derivative correction using the
order-0 Euler--Lagrange equation.  Cheyette~\cite{Cheyette1988} adapts
this to effective Standard Model Lagrangians; Simon~\cite{Simon1990}
formalises the perturbative ghost-suppression argument; Eliezer \&
Woodard~\cite{EliezerWoodard1989} prove the ghost branch contains
``runaway pseudosolutions'' that are excised by the iterative reduction;
Bender \& Mannheim~\cite{BenderMannheim2008} apply the framework to
the Pais--Uhlenbeck oscillator.  All assume the $\varepsilon = 0$ theory
is non-degenerate (every field propagates at the leading order).
\textbf{They fail on constraint promotion} because the substitution
rule $\partial_t^2 h_c \to \dots$ does not exist when $h_c$ is constrained
at order zero.  TIDAL's v6 chose Parker--Simon~\cite{ParkerSimon1993}
iterative source-evaluation \emph{precisely} to avoid this failure mode
on the EOM side; a Lagrangian analogue is not yet worked out.

\paragraph{EFT field redefinitions.}
Burgess \& Williams~\cite{BurgessWilliams2014} and the modern systematic
treatment of Solomon \& Trodden~\cite{SolomonTrodden2018} (see also
Manohar's lectures~\cite{Manohar2018}) use field redefinitions
$q \to q + \varepsilon\,\delta q[q, \partial q]$ to eliminate higher-derivative
operators perturbatively.  The redefinition is constructed so that the
shift in the Lagrangian cancels the unwanted derivative terms.
Mathematically this is equivalent to the JLM substitution and shares the
same architectural assumption: the leading-order kinetic matrix must be
invertible.  Glavan~\cite{Glavan2018} compares JLM, field redefinitions,
and Cayuso--Lehner side-by-side; all three break down at constraint
promotion.

\paragraph{Cayuso--Lehner ``fix the equations''.}
Cayuso \& Lehner~\cite{CayusoLehner2017} propose a non-iterative numerical
scheme that evolves the unreduced higher-derivative system with a
parameter-dependent damping term.  The method is intended for the
\emph{strong-coupling} regime where $\varepsilon$ is not small.  It does
not produce a perturbative Hamiltonian for our case, and the damping
parameter introduces theory-dependent tuning rather than an order-by-order
expansion.

\paragraph{Faddeev--Jackiw symplectic reduction.}
Faddeev \& Jackiw~\cite{FaddeevJackiw1988} provide an alternative to
Dirac's constraint analysis that operates directly on the symplectic
2-form.  Iterative bordering algorithms~\cite{DiazSaldivia2025} extend
this to systems with multiple stages of constraints.  Faddeev--Jackiw
is computationally more efficient than Dirac for second-class systems,
\emph{but does not avoid the Ostrogradsky degree-of-freedom count}.
For the constraint-promotion case, the Faddeev--Jackiw 2-form has rank
deficiency at $\varepsilon = 0$ that is lifted at $\varepsilon \ne 0$ --- exactly
the rank-jumping behaviour responsible for the discontinuity.  No published
Faddeev--Jackiw treatment perturbs in~$\varepsilon$ across this rank jump.

\paragraph{Cubic-PGT redesign (Bahamonde \emph{et al.}~2024--2025).}
The recent program~\cite{bahamonde2025cubic,bahamonde2025torsion}
adds cubic curvature invariants to the quadratic-PGT action and
fine-tunes the joint coefficients to remove the ghost
non-perturbatively.  This is a redesign of the theory --- new dynamical
content is added, with algebraic relations among the quadratic and
cubic coefficients enforcing ghost-freedom.  For an experimental or
phenomenological setting that asks ``what does the $b_5\,\tilde{R}^2$
theory predict, perturbatively in $b_5$?'' this approach is not
applicable; it answers a different question (``what is a related
ghost-free PGT?'').

\paragraph{DHOST degeneracy condition.}
Beyond-Horndeski / DHOST theories~\cite{Langlois2018} are constructed by
imposing a degeneracy condition on the kinetic matrix that ensures the
Ostrogradsky mode is removed at the Lagrangian level by a primary
constraint.  This is constraint-by-design: a class of theories \emph{built}
to be free of ghosts.  It does not provide a perturbative reduction
method for a pre-existing Lagrangian.

\paragraph{Stelle quadratic gravity.}
The classical reference for the canonical Hamiltonian of quadratic
gravity is Stelle~1977~\cite{Stelle1977}.  This treatment retains the
Ostrogradsky ghost as a real (massive) degree of freedom and is
\emph{non-perturbative} in the curvature-squared coupling.  It is the
``honest'' Hamiltonian for the full quadratic theory but does not give
a perturbative Hamiltonian valid only at small coupling.

\paragraph{Conclusion of the survey.}
No published method handles the case where $\varepsilon$ continuously
promotes an algebraic constraint to a higher-derivative dynamical mode.
The closest acknowledged obstruction is Henneaux \&
Teitelboim's~\cite{HenneauxTeitelboim1992} ``irregular constraint
surface'', identified in canonical analysis but not bridged by any
perturbative algorithm.

\subsection{TIDAL's resolution and current behaviour}
\label{sec:pr-cb-resolution}

In light of the architectural barrier, TIDAL's v6 Phase~2 LPS is
deliberately scoped to the cases where it is mathematically justified
and aborts \emph{loudly} on detection of the constraint-promotion case.

\paragraph{Detection.}
LPS examines the input Lagrangian for components carrying
$\partial_t^2$ derivatives (\texttt{DetectMaxTimeOrder} in
\texttt{PerturbativeReduction.wl}).  For each such component~$h_c$, LPS
inspects the order-0 EOM: if the EOM contains $\partial_t^2 h_c$, $h_c$
is dynamical at the leading order and substitution proceeds; if not,
$h_c$ is constraint-promoted, and LPS aborts via Wolfram's
\texttt{Throw} mechanism.

\paragraph{Error behaviour.}
The Throw is caught at the WLS-template level (commit \texttt{742c01a},
\texttt{tidal/cli/\_derive.py}); on detection the script prints a
clear diagnostic naming the constraint-promoted fields and exits with
status~1.  Python's lenient handling of post-derivation JSON is
\emph{disabled} by an mtime check: if the on-disk JSON was not modified
during the failed run, the lenience is skipped and the failure
propagates.  This prevents stale JSONs from masking real errors --- a
defensive measure motivated by the principle that silent miscorrection
is worse than loud failure, and aligned with the project policy of
fail-loud over fall-back.

\paragraph{Affected examples.}
Three TIDAL examples trigger the abort:
\texttt{graviton\_torsion}, \texttt{torsion\_gertsenshtein}, and
\texttt{torsion\_gertsenshtein\_combined}.  All carry $b_5\,\tilde{R}^2$
or analogous quadratic-curvature/torsion couplings.  Their
pre-Phase~2 JSONs (containing residual \texttt{mixed\_2\_*} operators in
the Hamiltonian section) remain on disk and continue to function for
measurement paths that re-evaluate the EOM at snapshot time (rather
than finite-differencing the velocity slot across snapshots --- a
known instability when snapshots are widely spaced).  The \emph{equation}
side (where Parker--Simon iterative reduction succeeds) is unaffected.

\paragraph{Unaffected examples.}
Theories without $b_5\,\tilde{R}^2$-style constraint-promotion proceed
normally.  Notably, Euler--Heisenberg
($\frac{\rho}{8}(F\cdot F)^2 + \frac{\sigma}{8}(\varepsilon \cdot F\cdot F)^2$)
has a Lagrangian that is naturally first-order in time after IBP, so
LPS is a no-op for that theory.  Pais--Uhlenbeck-style perturbations
to dynamical Lagrangians proceed cleanly --- this is the textbook case
that the substitution-based perturbative-reduction literature was
designed for, and TIDAL's LPS handles it correctly (covered by
\texttt{tests/wolfram/test\_lps.wls}).

\subsection{Future work}
\label{sec:pr-cb-future}

The constraint-promotion case is open research, not engineering.  Issue
\#321 in TIDAL's tracker outlines three directions, none of which has a
published implementation for this specific problem.

\paragraph{(a) Algebraic-constraint resolution with chain-rule expansion.}
Solve the order-0 algebraic constraint for $h_c$ in terms of dynamical
fields, chain-differentiate twice in time, and substitute back into the
Lagrangian.  This is what was attempted during Phase~2 development
(commit branch reverted before merge): the procedure does converge
mathematically, but the chain rule introduces $\partial_t^2$ of
\emph{dynamical} fields, which then require their own substitution.
On TIDAL's largest constraint-promotion theory
(\texttt{torsion\_gertsenshtein}, 38 fields), the substitution failed
to complete within reasonable resources (5+~min, 2+~GB) and
correctness was uncertain (rule-chain convergence and Series-truncation
interaction not formally verified).  A successful implementation would
need both algorithmic optimisation \emph{and} a correctness proof; both
are research-level.

\paragraph{(b) Parker--Simon analogue at the Lagrangian level.}
The v6 EOM-side success rests on Parker--Simon's iterative
\emph{source} treatment of higher-derivative corrections: the
correction term acts as a Duhamel source on the second-order base
operator, never appearing in the LHS.  No published analogue exists for
the Lagrangian, because the Lagrangian is a static functional with no
``evolution'' to convert into a Duhamel kernel.  A novel theoretical
construction would be needed.

\paragraph{(c) Hamiltonian directly from the v6-reduced EOM.}
Given the v6 Parker--Simon-reduced equations of motion, attempt to
reconstruct an effective Hamiltonian by antiderivative integration:
$H_0 = \tfrac{1}{2}\,\pi\, M_0^{-1}\, \pi + V_0(q)$ with
$V_0(q) = -\int K_0(q)\,dq$, plus an order-1 correction
$\varepsilon\, V_1(q)$ from the Pass-1 source.  The
antiderivative-reconstruction step is fragile for gauge-coupled
multi-field systems (no uniqueness guarantee, and the Pass-1 source
$K_1 - M_1 M_0^{-1} K_0$ is not in general the gradient of a scalar
potential).  This direction would need a constructive uniqueness or
gauge-invariant form theorem before being implementable.

\paragraph{Symplectic-perturbation lifting.}
A research direction not in the TIDAL issue list but suggested by the
Faddeev--Jackiw analysis: develop a perturbation theory of the
symplectic 2-form across rank jumps.  Standard deformation
quantisation~\cite[for example]{HenneauxTeitelboim1992} treats
deformations of \emph{non-degenerate} symplectic structures; the
rank-jumping case appears in the symplectic-foliation literature but not
as an explicit perturbation theory.  This is an unpublished research
direction.

\paragraph{What ``solving'' the barrier looks like.}
A successful perturbative-Hamiltonian reduction for the constraint-promotion
case would produce, given $\mathcal{L}(q; \varepsilon) = \mathcal{L}_0(q) +
\varepsilon\, \mathcal{L}_1(q, \partial^2 q, \dots)$ with $\varepsilon$
promoting~$h_c$, a Hamiltonian
$H(q, \pi; \varepsilon) = H_0(q, \pi) + \varepsilon\, H_1(q, \pi)$ such
that:

\begin{enumerate}
  \item[(i)] $H$ depends only on $(q, \pi)$, with no Ostrogradsky higher
        momenta.
  \item[(ii)] Hamilton's equations from~$H$ reproduce, to order~$\varepsilon$,
        the Parker--Simon-reduced EOM of the v6 path.
  \item[(iii)] The reduction is well-defined for the multi-field
        gauge-theory case relevant to PGT torsion.
  \item[(iv)] The construction is theory-agnostic, not requiring per-theory
        fine-tuning.
\end{enumerate}

No published method achieves all four for the constraint-promotion case.

\subsection{Conclusion}
\label{sec:pr-cb-conclusion}

The constraint-promotion case is a structural barrier to perturbative
reduction of higher-derivative Hamiltonians, structurally analogous to
the vDVZ discontinuity and the Maxwell--Proca constraint reclassification.
Every published reduction method --- substitution, field redefinition,
fixing-the-equations, Faddeev--Jackiw, cubic-PGT redesign, DHOST --- shares
an architectural assumption that excludes this case.  TIDAL's v6 Phase~2
LPS therefore aborts loudly when constraint promotion is detected,
preserving correctness over partial output.  The affected
$b_5\,\tilde{R}^2$ PGT theories continue to function via their
pre-Phase-2 JSONs and the v6 EOM-side Parker--Simon path; the
perturbative Hamiltonian for these theories remains an open research
question.

This document records the architectural lesson and the relevant
literature so the lesson is not re-discovered, and provides three
research directions for future work.  Implementation would require
novel theoretical results, not engineering effort, and is therefore
deferred indefinitely with the understanding that the current TIDAL
behaviour is the architecturally honest answer.
